\section{Method}
\label{sec:method}

\subsection{Fixed Retrieval Packet}

Text retrieval combines BGE-M3 dense sentence similarity with BM25 using reciprocal-rank fusion (RRF) \cite{cormack2009rrf}. A pinned BGE reranker scores the fused text candidates. Image retrieval combines SigLIP2 image-text similarity with dense similarity between the claim and normalized image filenames, again using RRF. The exact model repositories and immutable revisions are recorded in the frozen configuration; the operational model cards are BGE-M3,\footnote{\url{https://huggingface.co/BAAI/bge-m3}} BGE-reranker-v2-m3,\footnote{\url{https://huggingface.co/BAAI/bge-reranker-v2-m3}} and SigLIP2.\footnote{\url{https://huggingface.co/google/siglip2-base-patch16-224}} The retained packet contains the highest-ranked 40 text and 20 image items. Every selector receives this identical packet.

\subsection{Membership Scorers}

For each claim-candidate pair, the stored feature vector contains the final score and its square, dense and lexical or filename scores, RRF and cross-encoder scores, reciprocal ranks, modality indicators, token-count features, available-candidate count, and duplicate-group size. Features are standardized using fitting-role statistics. A logistic scorer then estimates
\begin{equation}
q_{qi}=\sigma\!\left(\mathbf{w}^{\top}
\frac{\mathbf{x}_{qi}-\boldsymbol{\mu}}{\boldsymbol{s}}+b\right),
\label{eq:membership-scorer}
\end{equation}
where $q_{qi}$ is a ranking score rather than an assumed operating probability. We compare three organizations: one shared text-image model, separate modality models, and a shared model using the complete feature vector including cross-encoder scores. Each organization is trained with balanced and unweighted binary cross-entropy plus $L_2$ regularization. Balanced loss is useful for ranking rare evidence, but its output reflects an altered effective class prior and cannot be treated as calibrated.

Let $z_{qi}=\operatorname{logit}(q_{qi})$. A calibration map $g_m$ converts it to the membership probability used by the decision policy,
\begin{equation}
p_{qi}=g_m(z_{qi})\approx P(c_{qi}\in G_q\mid \mathbf{x}_{qi}),
\label{eq:calibration-map}
\end{equation}
with a separate fitted map for each modality. Six maps are compared:
\begin{enumerate}
    \item identity, which leaves the raw sigmoid probability unchanged;
    \item temperature scaling, $\sigma(z/T)$;
    \item prior shift, which fits an intercept correction;
    \item Platt mapping, $\sigma(az+b)$;
    \item beta calibration, which is linear in $\log p$ and $\log(1-p)$; and
    \item isotonic regression, a monotone nonparametric map.
\end{enumerate}
Calibration parameters are fitted on the calibration role. The scorer and map are chosen by average precision and Brier score on the separate model-selection role. We report Brier score, negative log likelihood (NLL), expected calibration error (ECE), mean predicted probability, empirical prevalence, and Brier skill relative to the prevalence-only predictor.

This separation prevents a calibrator from being selected on the locked-validation outcome. Average precision screens the scorer's ordering, while Brier score selects the probability map used by threshold and expected-utility decisions. The distinction is important: any monotone map preserves ranking but may substantially change the selected set.

\subsection{Evidence-Set Policies}

Figure~\ref{fig:policy-decoding} isolates the only changing operation: how one shared calibrated ranking is converted into an evidence set.

\begin{figure}[t]
    \centering
    \resizebox{\columnwidth}{!}{\begin{tikzpicture}[
  font=\scriptsize,
  policy/.style={draw=accentblue, rounded corners=2pt, fill=softblue,
    minimum height=9mm, text width=18mm, align=center, inner sep=2.5pt},
  rule/.style={draw=black!45, rounded corners=2pt, fill=softgray,
    minimum height=9mm, text width=29mm, align=center, inner sep=2.5pt},
  selected/.style={draw=accentteal, rounded corners=2pt, fill=accentteal!9,
    minimum height=9mm, text width=25mm, align=center, inner sep=2.5pt},
  shared/.style={draw=accentorange, rounded corners=2pt, fill=accentorange!10,
    minimum height=10mm, text width=75mm, align=center, inner sep=3pt},
  note/.style={draw=black!40, dashed, rounded corners=2pt,
    text width=75mm, align=center, inner sep=2.5pt, text=black!70},
  arrow/.style={-{Latex[length=1.8mm]}, thick, draw=black!60}
]

\node[shared] (ranking) {Shared calibrated ranking\\
$c_{q1},\ldots,c_{q60}$ with $p_{q1}\geq\cdots\geq p_{q60}$};

\node[policy, below=4mm of ranking, xshift=-29mm] (fixed) {Fixed top-$k$};
\node[rule, right=2.5mm of fixed] (fixedrule) {Use tuned global size $k=5$};
\node[selected, right=2.5mm of fixedrule] (fixedset) {Return the first 5 items};

\node[policy, below=2.5mm of fixed] (threshold) {Calibrated threshold};
\node[rule, right=2.5mm of threshold] (thresholdrule) {Keep items with $p_{qi}\geq\tau$};
\node[selected, right=2.5mm of thresholdrule] (thresholdset) {Prefix ending at the last score above $\tau$};

\node[policy, below=2.5mm of threshold] (count) {Learned count};
\node[rule, right=2.5mm of count] (countrule) {Predict set size $\widehat{k}_q$};
\node[selected, right=2.5mm of countrule] (countset) {Return the first $\widehat{k}_q$ items};

\node[policy, below=2.5mm of count] (expected) {Expected-F1};
\node[rule, right=2.5mm of expected] (expectedrule) {Evaluate $U_q(k)$ for $k=0{:}60$};
\node[selected, right=2.5mm of expectedrule] (expectedset) {Return the first $k_q^*$ items; break ties short};

\draw[arrow] (fixed) -- (fixedrule);
\draw[arrow] (fixedrule) -- (fixedset);
\draw[arrow] (threshold) -- (thresholdrule);
\draw[arrow] (thresholdrule) -- (thresholdset);
\draw[arrow] (count) -- (countrule);
\draw[arrow] (countrule) -- (countset);
\draw[arrow] (expected) -- (expectedrule);
\draw[arrow] (expectedrule) -- (expectedset);

\node[note, below=4mm of expectedrule, xshift=0mm] (scope)
{Only the set decision changes; candidate identities, ordering, and calibrated probabilities are shared.};

\coordinate (spinetop) at ([xshift=-4mm]fixed.west);
\coordinate (spinebottom) at ([xshift=-4mm]expected.west);
\draw[thick, draw=black!60] (ranking.south) -- ++(0,-1.5mm) -| (spinetop);
\draw[thick, draw=black!60] (spinetop) -- (spinebottom);
\draw[arrow] (spinetop) -- (fixed.west);
\draw[arrow] (spinetop |- threshold.west) -- (threshold.west);
\draw[arrow] (spinetop |- count.west) -- (count.west);
\draw[arrow] (spinebottom) -- (expected.west);
\end{tikzpicture}
}
    \caption{Decision policies applied to the same calibrated ranking. Fixed top-5 uses one global size; thresholding uses an absolute cutoff; learned count predicts the size; expected-F1 evaluates every prefix and returns the shortest utility-maximizing one. The displayed values illustrate the procedure; they are not measurements.}
    \label{fig:policy-decoding}
\end{figure}

\paragraph{Fixed top-$k$.}
We tune $k\in\{1,2,3,5,10,20,40,60\}$ and select the top $k$ candidates. The resulting control is top-5.

\paragraph{Calibrated threshold.}
For threshold $\tau$, the selector returns every item with $p_{qi}\geq\tau$, subject to the common packet limit. The threshold is selected from $\{0.1,\ldots,0.9\}$.

\paragraph{Learned count.}
A count model predicts the number of evidence items, after which the corresponding probability-ranked prefix is returned.

\paragraph{Expected-F1 decoding.}
Candidates are sorted by calibrated probability. For each prefix length $k$, a dynamic program computes the exact expected F1 under independent Bernoulli membership variables. The selected size is
\begin{equation}
k_q^*=\arg\max_{0\leq k\leq 60}
\mathbb{E}\!\left[F_1(S_{q,k},Y_q)\mid p_{q1:60}\right].
\label{eq:expected-f1}
\end{equation}
For prefix $S_{q,k}$, the utility can be written
\begin{equation}
U_q(k)=\mathbb{E}\!\left[
\frac{2T_{q,k}}{k+T_{q,k}+R_{q,k}}
\;\middle|\;p_{q1:60}\right],
\label{eq:expected-f1-utility}
\end{equation}
where $T_{q,k}$ is the number of relevant items inside the prefix and $R_{q,k}$ is the number outside it. Two Poisson-binomial recurrences compute their distributions and the expectation exactly. Ties retain the shorter prefix. This rule uses the full membership distribution rather than a single cutoff or a separately predicted count.

We also test two extensions. The first subtracts $\lambda k$ from expected utility. The second additionally subtracts $\gamma R(S_{q,k})$, where $R$ counts same-source redundancy. Size penalties span 0 to 0.10 and redundancy penalties span 0 to 0.005. These terms are retained as ablations because the fitted optima reveal whether extra decoder structure contributes beyond calibration.

\subsection{Decision Scope and Assumptions}

The comparison changes only the set decision. Candidate identities, ranks, feature values, scorer outputs, and validation claims are shared across policies. Expected-F1 assumes conditionally independent Bernoulli memberships and is restricted to prefixes of the frozen ranking; it neither reorders candidates nor retrieves a missing item. The released complete-gold mapping is treated as the evaluation target, although Sec.~\ref{sec:limitations} discusses possible annotation incompleteness. These boundaries distinguish a decoder error from a retrieval error and prevent an adaptive policy from receiving a more favorable candidate universe.

\subsection{Shared Verifier}

Evidence policies can improve set overlap while harming the downstream verdict. To isolate that possibility, one verifier architecture and training procedure is shared across policies. Only the selected evidence changes. We compare verdict macro-F1 on the same 1,490 claims and use a declared non-inferiority margin of $-0.01$ for expected-F1 relative to fixed top-5.
